\subsection{Finalizing provisioning configuration} \label{sec:assemble_bootstrap}

\Warewulf{} provisions a node with an image then customizes it with overlays.
This section highlights creation of the node image and overlays, followed by the
registration of desired compute nodes.

\subsubsection{Two-Stage Provisioning (optional)}

\Warewulf{} optionally supports two-stage provisioning to ramfs or disk on the
compute nodes. The nodes are still stateless, but the image is copied to disk
during the provisioning process and is reprovisioned on every boot. This can be
useful for cases when the compute node images are large, for example GPU
drivers, or the compute nodes have limited memory.

To configure two-stage provisioning, install Dracut in the compute node image.
The example below provisions to ramfs, which is the default. To provision to
disk, see the next section.

% begin_ohpc_run
% ohpc_command if [[ ${enable_dracut} -eq 1 ]];then
% ohpc_comment_header Two-Stage provision to ramfs
\begin{lstlisting}[language=bash,literate={-}{-}1,keywords={},upquote=true,literate={BOSVER}{\baseos{}}1]
## Add rd.shell to kernel arguments in the "default" profile.
[sms](*\#*) yq -y -i '.nodeprofiles.default.kernel.args += ["rd.shell"]' \
  /etc/warewulf/nodes.conf

## Install Dracut in the image
[sms](*\#*) wwctl image exec --build=false BOSVER -- \
  /usr/bin/apt -y install warewulf-dracut-ohpc

## Configure Dracut. Note the leading and trailing spaces in the quotes in "add_dracutmodules"
[sms](*\#*) echo 'hostonly="no"' > $CHROOT/etc/dracut.conf.d/wwinit.conf
[sms](*\#*) echo 'add_dracutmodules+=" wwinit "' >> $CHROOT/etc/dracut.conf.d/wwinit.conf

## Enable Dracut boot for compute nodes that use the "nodes" profile
[sms](*\#*) wwctl profile set --yes nodes --tagadd IPXEMenuEntry=dracut
\end{lstlisting}
% ohpc_command fi
% end_ohpc_run

Optionally, configure the compute node to provision to disk. Two stage
provisioning (previous step) must also be enabled. The compute node will use,
\textbf{and erase}, the disk specified by \texttt{node\_disk}.

% begin_ohpc_run
% ohpc_command if [[ ${enable_dracut} -eq 1 && ${enable_todisk} -eq 1 ]];then
% ohpc_comment_header Configure two-stage provision-to-disk
\begin{lstlisting}[language=bash,literate={-}{-}1,keywords={},upquote=true,literate={BOSVER}{\baseos{}}1]
## Install disk tools
[sms](*\#*) wwctl image exec --build=false BOSVER -- /usr/bin/apt -y install ignition gdisk

## Add ignition Dracut module.
[sms](*\#*) echo 'add_dracutmodules+=" ignition "' >> $CHROOT/etc/dracut.conf.d/wwinit.conf

## Create the target "rootfs" partition and filesystem
[sms](*\#*) wwctl profile set --yes nodes \
  --diskname ${node_disk} --diskwipe \
  --partname rootfs --partcreate --partnumber 1 \
  --fsname rootfs --fswipe --fsformat ext4 --fspath /

## Enable provision-to-disk for compute nodes that use the "nodes" profile
[sms](*\#*) wwctl profile set nodes --yes --root=/dev/disk/by-partlabel/rootfs
\end{lstlisting}
% ohpc_command fi
% end_ohpc_run

\subsubsection{Build image image and overlays}

The bootstrap image includes the runtime kernel and associated modules, as well
as some simple scripts to complete the provisioning process. We explicitly rebuild
the image at the end because rebuilding the image is disabled in earlier steps.
It is good practice to rebuild the overlays when after configuration changes.

% begin_ohpc_run
% ohpc_comment_header Assemble bootstrap image \ref{sec:assemble_bootstrap}
\begin{lstlisting}[language=bash,literate={-}{-}1,keywords={},upquote=true,literate={BOSVER}{\baseos{}}1]
# Build Dracut initramfs and image
[sms](*\#*) wwctl image exec --build=false BOSVER -- /usr/bin/dracut --force --regenerate-all
[sms](*\#*) wwctl image build BOSVER
[sms](*\#*) wwctl overlay build
\end{lstlisting}
% end_ohpc_run

\subsubsection{Register nodes for provisioning}

Nodes can be registered for provisioning using the following syntax.
